\section{Introduction}
\label{sec:introduction}

A formulaic alpha is a short program that maps past prices and volumes of a
cross-section of assets to a score predicting next-period relative returns.
Public catalogues of such programs~\citep{kakushadze2016alphas} and decades of
anomaly research show both that the space is rich and that it is easy to
overfit. \citet{harvey2016cross} argue that the conventional $t>2$ hurdle is
far too permissive once the number of tested factors is acknowledged, and
backtest overfitting is a well-documented failure mode in quantitative
finance~\citep{bailey2014deflated,bailey2017pbo}. Automated search worsens
this problem, because a search procedure can generate and test candidates
faster than any human analyst. Genetic
programming~\citep{koza1992genetic} and reinforcement
learning~\citep{yu2023alphagen} have both been used to mine formulaic alphas.

Large language models add a new kind of proposer. An LLM can write
syntactically valid expressions, explain them in economic language, and
refine them in response to feedback. LLM-guided program search has produced
verifiable discoveries in mathematics~\citep{romeraparedes2024funsearch}, and
interactive LLM systems for alpha mining have been
proposed~\citep{wang2023alphagpt}. Finance differs from mathematics in two ways
that matter for evaluation.
\begin{enumerate}
  \item \emph{Ground truth.} There is no verifier that certifies a factor. The
    only evidence is statistical, drawn from one non-repeatable history, so a
    method that proposes more candidates is guaranteed to find more spurious
    ones unless every trial is counted.
  \item \emph{Look-ahead through pre-training.} An LLM may have read which
    signals worked in which period. LLMs have been used to predict returns from
    text~\citep{lopezlira2023chatgpt}, and such predictions can suffer from
    look-ahead bias inherited from the training
    data~\citep{glasserman2023lookahead}. An apparent edge may therefore
    reflect memory rather than search.
\end{enumerate}

\paragraph{Research questions.} This paper asks four questions:
\begin{itemize}
  \item \textbf{RQ1.} At an equal trial budget, does LLM-guided search yield
    frozen factor sets that perform better on a sealed test window than random
    grammar sampling and genetic programming, after multiplicity control that
    counts all trials and confirms on data no proposer has seen?
  \item \textbf{RQ2.} How do the methods differ in search efficiency?
  \item \textbf{RQ3.} How much of any LLM advantage is explained by
    rediscovering textbook factors, or by period-specific knowledge absorbed
    during pre-training?
  \item \textbf{RQ4.} Which components of the LLM proposer matter?
\end{itemize}
The hypotheses and decision rules were written before any LLM or real-data
experiment, but after the first baseline runs on the synthetic benchmark: a
review of those runs led us to revise the benchmark's planted signals and the
synthetic hypothesis, and we disclose this. The rules are given in the
project's research plan, whose hash is to be published before the first LLM
call, and are summarized in Section~\ref{sec:setup}.

\paragraph{Approach.} We treat the proposer as an untrusted component and
build the evaluation protocol around it (Section~\ref{sec:method}):
\begin{itemize}
  \item \emph{The language.} Expressions are written in a typed language
    whose operators only look backward, so no valid expression can read
    future data.
  \item \emph{Search.} The search loop evaluates on a price panel truncated
    at the end of the formation window. The proposer receives only
    formation-window statistics, through closed types whose only free-text
    fields hold the proposal's own text and a validator or fixed harness
    message.
  \item \emph{The ledger.} Every processed proposal is written to a
    hash-chained ledger, including those that fail validation. A formation
    screen counts every budgeted trial and merges behaviourally identical
    ones; because adaptive proposers see formation results, the claim rests
    on a second, confirmatory test on the validation window.
  \item \emph{The hold-out.} The selected set is committed by hash and
    revealed on the test window at most once; a project-wide registry fixes
    each real-data study's windows in advance, counts exploratory
    evaluations and caps reveals.
  \item \emph{Contamination.} Diagnostics compare test performance before and
    after the model's documented knowledge cutoff, and measure distance from
    a library of classic factors on factor \emph{values}.
\end{itemize}

\paragraph{Contributions.}
\begin{enumerate}
  \item An auditable protocol and open implementation for LLM-driven factor
    search. It covers the language, the information barrier, the trial
    ledger, the two-stage multiplicity control, the commit-then-reveal
    hold-out and the study registry, and every control is covered by offline
    tests.
  \item A planted-alpha synthetic benchmark with null markets, whose ground
    truth mixes textbook signals, a documented factor family and a signal
    drawn at random by a procedure fixed in code before the draw was run
    (self-attested, not externally registered); only the drawn signal
    speaks to search beyond prior knowledge. Results on it are reported for
    two baselines.
  \item \todo{An equal-budget comparison of LLM, random-grammar and
    genetic-programming search on US equities (pending).}
  \item \todo{Contamination analysis: knowledge-cutoff difference-in-differences,
    novelty stratification and anonymization probes (pending).}
\end{enumerate}

\paragraph{Current evidence.} The only experiment completed so far is the
harness validation on synthetic data (Section~\ref{sec:results-synthetic}).
With \SynBudget{} trials per run, neither baseline selected a factor in any
of \SynNullRuns{} null-market runs, although genetic programming's formation
screen alone passed spurious expressions. On planted markets both baselines
recovered only textbook signals and never found the drawn signal, and
genetic programming was not clearly better than random sampling. This leaves
room to test whether LLM guidance helps. \todo{State the LLM and real-data findings once the
pre-registered experiments have been run; do not edit this paragraph before
then.}
